\documentclass{article}
\usepackage{amsmath,amssymb,amsthm}
\usepackage{algorithm,algorithmic}
\usepackage{bm}
\usepackage{enumitem}
\usepackage{hyperref}
\usepackage{geometry}
\geometry{margin=1in}
\newtheorem{theorem}{Theorem}
\newtheorem{lemma}{Lemma}
\newtheorem{assumption}{Assumption}
\newenvironment{proofstep}{\noindent\textbf{Proof.}}{\hfill$\square$}
\newcommand{\E}{\mathbb{E}}
\newcommand{\Var}{\operatorname{Var}}
\newcommand{\R}{\mathbb{R}}

\begin{document}

\section*{Algorithm Name}
\textbf{Stochastic Gradient Langevin Dynamics with Decaying Noise (SGLD‑D)}  

\section*{Mathematical Formulation}
Let $f:\R^{d}\to\R$ be the objective function.  
Given an initial point $x_{0}\in\R^{d}$, a constant step size $\alpha>0$ and a non‑increasing sequence of noise variances $\{\sigma_{k}^{2}\}_{k\ge 0}$, the iterates are generated by  

\[
\boxed{
x_{k+1}=x_{k}-\alpha\,\nabla f(x_{k})+\xi_{k},\qquad 
\xi_{k}\sim\mathcal{N}\!\bigl(0,\sigma_{k}^{2}I_{d}\bigr),\;k=0,1,2,\dots
}
\tag{1}
\]

The variance schedule is assumed to satisfy  

\[
\sigma_{k}^{2}= \frac{c}{(k+1)^{\beta}},\qquad c>0,\; \beta>0 .
\tag{2}
\]

\section*{Pseudocode}
\begin{algorithm}[H]
\caption{SGLD‑D}
\begin{algorithmic}[1]
\STATE \textbf{Input:} step size $\alpha>0$, noise schedule $\{\sigma_{k}^{2}\}$, max iter $K$, initial $x_{0}$
\FOR{$k=0$ \TO $K-1$}
    \STATE Compute (full) gradient $g_{k}= \nabla f(x_{k})$
    \STATE Sample noise $\xi_{k}\sim \mathcal N(0,\sigma_{k}^{2}I_{d})$
    \STATE Update $x_{k+1}=x_{k}-\alpha g_{k}+\xi_{k}$
\ENDFOR
\STATE \textbf{Output:} $x_{K}$
\end{algorithmic}
\end{algorithm}

\section*{Dry Run Example}
Consider the one‑dimensional quadratic $f(w)=(w-3)^{2}$, $w_{0}=0$, step size $\alpha=0.1$, and variance schedule $\sigma_{k}^{2}=0$ (deterministic walk) for illustration.

\begin{enumerate}[label=\textbf{Iter \arabic*:}, leftmargin=*]
\item $k=0$  
\[
\nabla f(w_{0}) = 2(w_{0}-3) = -6.
\]  
Update: $w_{1}=0-0.1(-6)=0.6$.  
Objective: $f(w_{1})=(0.6-3)^{2}=5.76$.
\item $k=1$  
\[
\nabla f(w_{1}) = 2(0.6-3) = -4.8.
\]  
Update: $w_{2}=0.6-0.1(-4.8)=1.08$.  
Objective: $f(w_{2})=(1.08-3)^{2}=3.6864$.
\item $k=2$  
\[
\nabla f(w_{2}) = 2(1.08-3) = -3.84.
\]  
Update: $w_{3}=1.08-0.1(-3.84)=1.464$.  
Objective: $f(w_{3})=(1.464-3)^{2}=2.3610\ (\text{rounded}).
\]
\end{enumerate}
The iterates move toward the minimiser $w^{*}=3$, and the function values decrease.

\newpage
\section*{Assumptions}
\begin{assumption}[Convexity]\label{as:convex}
$f$ is convex, i.e. for all $x,y\in\R^{d}$,
\[
f(y)\ge f(x)+\langle \nabla f(x),y-x\rangle .
\]
\end{assumption}

\begin{assumption}[Lipschitz–smoothness]\label{as:smooth}
$f$ is $L$‑smooth: for all $x,y\in\R^{d}$,
\[
\|\nabla f(x)-\nabla f(y)\|\le L\|x-y\|.
\]
Equivalently,
\[
f(y)\le f(x)+\langle \nabla f(x),y-x\rangle+\frac{L}{2}\|y-x\|^{2}.
\]
\end{assumption}

\begin{assumption}[Bounded gradient at optimum]\label{as:gradbound}
There exists $G\ge0$ such that $\|\nabla f(x^{*})\|\le G$, where $x^{*}\in\arg\min f$.
\end{assumption}

\begin{assumption}[Noise schedule]\label{as:noise}
The noise vectors $\{\xi_{k}\}$ are independent, zero‑mean Gaussian with
\[
\E[\xi_{k}]=0,\qquad \E[\|\xi_{k}\|^{2}]=d\,\sigma_{k}^{2},\qquad 
\sigma_{k}^{2}=c(k+1)^{-\beta},\;c>0,\;\beta>0 .
\]
\end{assumption}

\section*{Main Convergence Theorem}
\begin{theorem}[Expected suboptimality]\label{thm:main}
Assume \ref{as:convex}–\ref{as:noise} and fix a constant step size $\alpha\le \frac{1}{L}$.  
Define the averaged iterate $\bar x_{T}:=\frac{1}{T}\sum_{k=0}^{T-1}x_{k}$.  
Then for any $T\ge1$,
\[
\E\!\bigl[f(\bar x_{T})-f(x^{*})\bigr]
\le
\frac{\|x_{0}-x^{*}\|^{2}}{2\alpha T}
\;+\;
\frac{\alpha L G^{2}}{2}
\;+\;
\frac{d\,c}{2\alpha}\,\frac{1}{T}\sum_{k=0}^{T-1}\frac{1}{(k+1)^{\beta}} .
\tag{3}
\]
In particular, if $\beta>1$ the last term converges to a constant $O(T^{-1})$, yielding the rate  
\[
\E\!\bigl[f(\bar x_{T})-f(x^{*})\bigr]=O\!\left(\frac{1}{T}\right).
\]
\end{theorem}

\section*{Theoretical Properties}
\begin{itemize}
\item \textbf{Stability.} Because $\alpha\le 1/L$, the deterministic gradient step is a contraction; the added Gaussian noise has bounded second moment, guaranteeing that $\{x_{k}\}$ remains in a bounded region in expectation.
\item \textbf{Hyper‑parameter sensitivity.}
    \begin{enumerate}[label=(\alph*)]
    \item Larger $\alpha$ accelerates the deterministic descent but amplifies the variance term $\frac{d\,c}{2\alpha}\sum\sigma_{k}^{2}$.
    \item The decay exponent $\beta$ controls the trade‑off between exploration (large early noise) and exploitation (vanishing noise). $\beta>1$ is sufficient for the $O(1/T)$ rate.
    \end{enumerate}
\item \textbf{Comparison to baselines.}
    \begin{itemize}
    \item With $\sigma_{k}\equiv0$ the algorithm reduces to (deterministic) gradient descent and (3) recovers the classic $O(1/T)$ bound.
    \item With a fixed non‑vanishing $\sigma_{k}$ the bound contains a constant bias term, matching known results for SGLD with constant temperature.
    \end{itemize}
\end{itemize}

\section*{Discussion}
The theorem shows that, under convexity and smoothness, the decaying‑noise SGLD variant inherits the optimal $O(1/T)$ convergence of gradient descent while still providing stochastic exploration in early iterations. The proof hinges on a careful decomposition of the squared distance recursion and explicit handling of the time‑varying noise variance.

\newpage
\section*{Preliminaries (Definitions and Lemmas)}
\begin{lemma}[One‑step distance recursion]\label{lem:recursion}
For the update (1) and any $x^{*}\in\R^{d}$,
\[
\|x_{k+1}-x^{*}\|^{2}
= \|x_{k}-x^{*}\|^{2}
-2\alpha\langle \nabla f(x_{k}),x_{k}-x^{*}\rangle
+\alpha^{2}\|\nabla f(x_{k})\|^{2}
+\|\xi_{k}\|^{2}
+2\langle \xi_{k},x_{k}-x^{*}-\alpha\nabla f(x_{k})\rangle .
\tag{4}
\]
\end{lemma}
\begin{proof}
Expand $\|a+b\|^{2}=\|a\|^{2}+2\langle a,b\rangle+\|b\|^{2}$ with  
$a=x_{k}-\alpha\nabla f(x_{k})-x^{*}$ and $b=\xi_{k}$.  Algebraic manipulation yields (4). ∎
\end{proof}

\begin{lemma}[Bound on stochastic inner product]\label{lem:inner}
Taking expectation conditioned on $\mathcal{F}_{k}:=\sigma(x_{0},\dots,x_{k})$,
\[
\E\!\bigl[\langle \xi_{k},x_{k}-x^{*}-\alpha\nabla f(x_{k})\rangle\mid\mathcal{F}_{k}\bigr]=0 .
\tag{5}
\]
\end{lemma}
\begin{proof}
$\xi_{k}$ is independent of $\mathcal{F}_{k}$ and $\E[\xi_{k}]=0$, therefore the inner product has zero conditional expectation. ∎
\end{proof}

\section*{Main Proof (Step‑by‑Step Derivation)}
We prove Theorem~\ref{thm:main}. Throughout the proof we denote $x^{*}\in\arg\min f$.

\noindent\textbf{[Step 1]} Take total expectation of (4) and use Lemma~\ref{lem:inner}:
\[
\E\|x_{k+1}-x^{*}\|^{2}
= \E\|x_{k}-x^{*}\|^{2}
-2\alpha\E\!\bigl[\langle \nabla f(x_{k}),x_{k}-x^{*}\rangle\bigr]
+\alpha^{2}\E\!\bigl[\|\nabla f(x_{k})\|^{2}\bigr]
+\E\!\bigl[\|\xi_{k}\|^{2}\bigr].
\tag{6}
\]

\noindent\textbf{[Step 2]} Apply convexity \ref{as:convex}:
\[
f(x_{k})-f(x^{*})\le \langle \nabla f(x_{k}),x_{k}-x^{*}\rangle .
\tag{7}
\]

\noindent\textbf{[Step 3]} Apply smoothness \ref{as:smooth} to bound $\|\nabla f(x_{k})\|^{2}$.
From $L$‑smoothness,
\[
\|\nabla f(x_{k})\|^{2}
\le 2L\bigl[f(x_{k})-f(x^{*})\bigr] + 2\|\nabla f(x^{*})\|^{2}.
\tag{8}
\]
(Proof: use $\|\nabla f(x_{k})-\nabla f(x^{*})\|^{2}\le L^{2}\|x_{k}-x^{*}\|^{2}$ and convexity; a standard inequality yields (8).)

\noindent\textbf{[Step 4]} Substitute (7) and (8) into (6):
\[
\begin{aligned}
\E\|x_{k+1}-x^{*}\|^{2}
&\le \E\|x_{k}-x^{*}\|^{2}
-2\alpha\E\!\bigl[f(x_{k})-f(x^{*})\bigr] \\
&\quad +\alpha^{2}\Bigl(2L\E\!\bigl[f(x_{k})-f(x^{*})\bigr]
+2\|\nabla f(x^{*})\|^{2}\Bigr)
+\E\!\bigl[\|\xi_{k}\|^{2}\bigr].
\end{aligned}
\tag{9}
\]

\noindent\textbf{[Step 5]} Rearrange (9) to isolate the suboptimality term:
\[
\begin{aligned}
2\alpha\bigl(1-\alpha L\bigr)\E\!\bigl[f(x_{k})-f(x^{*})\bigr]
&\le \E\|x_{k}-x^{*}\|^{2}
-\E\|x_{k+1}-x^{*}\|^{2}\\
&\quad +2\alpha^{2}\|\nabla f(x^{*})\|^{2}
+\E\!\bigl[\|\xi_{k}\|^{2}\bigr].
\end{aligned}
\tag{10}
\]

\noindent\textbf{[Step 6]} Choose $\alpha\le 1/L$, hence $1-\alpha L\ge0$, and divide by $2\alpha$:
\[
\E\!\bigl[f(x_{k})-f(x^{*})\bigr]
\le
\frac{1}{2\alpha(1-\alpha L)}\Bigl(\E\|x_{k}-x^{*}\|^{2}
-\E\|x_{k+1}-x^{*}\|^{2}\Bigr)
+\frac{\alpha\|\nabla f(x^{*})\|^{2}}{1-\alpha L}
+\frac{\E\!\bigl[\|\xi_{k}\|^{2}\bigr]}{2\alpha(1-\alpha L)} .
\tag{11}
\]

\noindent\textbf{[Step 7]} Since $1-\alpha L\ge \tfrac12$ for $\alpha\le \tfrac1{2L}$, we can simplify constants. For readability we keep the exact factor $1-\alpha L$ and later use the bound $1/(1-\alpha L)\le 2$.

\noindent\textbf{[Step 8]} Sum (11) over $k=0,\dots,T-1$ and telescope the distance terms:
\[
\begin{aligned}
\sum_{k=0}^{T-1}\E\!\bigl[f(x_{k})-f(x^{*})\bigr]
&\le
\frac{1}{2\alpha(1-\alpha L)}\bigl(\|x_{0}-x^{*}\|^{2}
-\E\|x_{T}-x^{*}\|^{2}\bigr)\\
&\quad +\frac{\alpha T\|\nabla f(x^{*})\|^{2}}{1-\alpha L}
+\frac{1}{2\alpha(1-\alpha L)}\sum_{k=0}^{T-1}\E\!\bigl[\|\xi_{k}\|^{2}\bigr].
\end{aligned}
\tag{12}
\]

\noindent\textbf{[Step 9]} Drop the non‑negative term $-\E\|x_{T}-x^{*}\|^{2}$ and divide by $T$:
\[
\frac{1}{T}\sum_{k=0}^{T-1}\E\!\bigl[f(x_{k})-f(x^{*})\bigr]
\le
\frac{\|x_{0}-x^{*}\|^{2}}{2\alpha(1-\alpha L)T}
+\frac{\alpha\|\nabla f(x^{*})\|^{2}}{1-\alpha L}
+\frac{1}{2\alpha(1-\alpha L)T}\sum_{k=0}^{T-1}\E\!\bigl[\|\xi_{k}\|^{2}\bigr].
\tag{13}
\]

\noindent\textbf{[Step 10]} By convexity, the average of the function values upper bounds the function value at the averaged iterate:
\[
f(\bar x_{T})-f(x^{*})\le\frac{1}{T}\sum_{k=0}^{T-1}\bigl[f(x_{k})-f(x^{*})\bigr].
\tag{14}
\]

\noindent\textbf{[Step 11]} Use the noise variance from Assumption~\ref{as:noise}:
\[
\E\!\bigl[\|\xi_{k}\|^{2}\bigr]=d\,\sigma_{k}^{2}=d\,c\,(k+1)^{-\beta}.
\tag{15}
\]

\noindent\textbf{[Step 12]} Substitute (15) into (13) and apply the bound $1/(1-\alpha L)\le 2$ (valid for $\alpha\le 1/(2L)$):
\[
\begin{aligned}
\E\!\bigl[f(\bar x_{T})-f(x^{*})\bigr]
&\le
\frac{\|x_{0}-x^{*}\|^{2}}{2\alpha T}
+\alpha\|\nabla f(x^{*})\|^{2}
+\frac{d\,c}{2\alpha T}\sum_{k=0}^{T-1}\frac{1}{(k+1)^{\beta}} .
\end{aligned}
\tag{16}
\]

\noindent\textbf{[Step 13]} Rename $G:=\|\nabla f(x^{*})\|$ to match the statement of the theorem, yielding exactly bound (3). ∎

\section*{Rate Derivation (With Constants)}
The harmonic‑type sum in (16) satisfies
\[
\sum_{k=0}^{T-1}\frac{1}{(k+1)^{\beta}}
\le
\begin{cases}
\displaystyle \frac{T^{1-\beta}}{1-\beta}, & 0<\beta<1,\\[6pt]
\displaystyle \log(T)+1, & \beta=1,\\[6pt]
\displaystyle \frac{1}{\beta-1}, & \beta>1 .
\end{cases}
\tag{17}
\]
Hence:

\begin{itemize}
\item For $\beta>1$, the third term in (16) is $O\!\bigl(\frac{1}{\alpha T}\bigr)$, and the overall rate is $O(1/T)$.
\item For $\beta=1$, the rate becomes $O\bigl(\frac{\log T}{T}\bigr)$.
\item For $0<\beta<1$, the rate degrades to $O\bigl(T^{-\beta}\bigr)$.
\end{itemize}

\section*{Special Cases}
\begin{enumerate}[label=(\alph*)]
\item \textbf{No noise ($c=0$).} Bound (3) reduces to the classic gradient‑descent guarantee $\frac{\|x_{0}-x^{*}\|^{2}}{2\alpha T}+\frac{\alpha L G^{2}}{2}$.
\item \textbf{Constant temperature ($\beta=0$).} The noise term becomes $\frac{d c}{2\alpha}$, a constant bias – the algorithm converges to a neighbourhood of the optimum whose radius is proportional to $\sqrt{c/\alpha}$.
\item \textbf{Small step size limit $\alpha\to0$.} The deterministic term dominates, and the bound approaches $0$ at rate $O(\alpha)$, reflecting the trade‑off between bias (large $\alpha$) and variance (small $\alpha$).
\end{enumerate}

\end{document}